\section*{Supporting information}

% This appendix provides supplementary details on the training configuration, human evaluation protocol, and additional experimental results that support the main text.

\subsection*{LoRA Training Hyperparameters}

Table~\ref{tab:hyperparams} lists the complete hyperparameter configuration used for doctor-specific LoRA training via LLaMA-Factory~\cite{zheng2024llamafactory}.

\begin{table}[H]
\caption{LoRA training hyperparameters.\label{tab:hyperparams}}
\begin{tabularx}{\textwidth}{lX}
\toprule
\textbf{Parameter} & \textbf{Value} \\
\midrule
Base model & LLaMA-3.2-1B/3B-Instruct, Qwen-2.5-7B-Instruct \\
LoRA rank $r$ & 8 \\
LoRA alpha $\alpha$ & 16 \\
LoRA dropout & 0.1 \\
Target modules & $\{q\_proj, k\_proj, v\_proj, o\_proj\}$ \\
Learning rate & $2 \times 10^{-4}$ \\
LR schedule & Cosine with warmup (10\% steps) \\
Batch size (per device) & 4 \\
Gradient accumulation steps & 8 \\
Effective batch size & 32 \\
Max epochs & 10 \\
Early stopping patience & 3 \\
Optimizer & AdamW ($\beta_1=0.9$, $\beta_2=0.999$) \\
Weight decay & 0.01 \\
Max sequence length & 2048 \\
Precision & BF16 mixed precision \\
Hardware & 8 $\times$ NVIDIA H20-3e (141GB) \\
\bottomrule
\end{tabularx}
\end{table}

\subsection*{LLM Judge Evaluation Protocol}

We adopt Qwen3.6-35B as an automated judge following the LLM-as-a-Judge paradigm~\cite{NEURIPS2023_91f18a12}. For each test query, the judge receives the generated response together with the target physician's department, specialty description, and three representative historical responses as context. The judge evaluates each response on three dimensions using 5-point Likert scales:

\begin{itemize}
\item \textbf{Personalization} (1--5): How well does the response match the target physician's characteristic communication style, terminology preferences, and clinical reasoning patterns?
\item \textbf{Clinical Accuracy} (1--5): Is the medical content correct, appropriate, and consistent with evidence-based practice for the presented complaint?
\item \textbf{Fluency} (1--5): Is the response grammatically correct, well-structured, and natural in its language use?
\end{itemize}

For preference assessment, the judge performs pairwise comparisons between all method pairs. Self-consistency is evaluated by repeating each judgment three times with different prompt orderings; the majority vote is taken as the final preference. The judge achieves a self-consistency rate of 94.2\%, indicating high reliability. 
% We note that Qwen3.6-35B belongs to the same model family as one of our evaluated base models (Qwen-2.5-7B-Instruct), which may introduce a self-preference bias. To mitigate this, we evaluate all methods---including the Qwen-2.5-7B-based variants---under identical judge configurations and report relative improvements rather than absolute scores as the primary comparison metric.

\subsection*{SCF Sensitivity Analysis}

Table~\ref{tab:scf-sensitivity} reports the sensitivity of SCF performance to the trim fraction $\tau$ with $K=120$. Performance is robust for $\tau \in [0.15, 0.30]$ and degrades gracefully as $K$ increases from 20 to 120, suggesting that the framework scales well to larger physician pools.

\begin{table}[H]
\caption{SCF sensitivity to trim fraction $\tau$ ($K=120$, 3B model). \label{tab:scf-sensitivity}}
\begin{tabularx}{\textwidth}{lCC}
\toprule
\textbf{Trim fraction $\tau$} & \textbf{Merged BLEU-4} &\textbf{Retention (\%)} \\
\midrule
0.10 & 9.3 & 89.4 \\
0.15 & 10.2 & 98.4 \\
0.20 & \textbf{10.4} & \textbf{100} \\
0.25 & 10.2 & 98.1 \\
0.30 & 9.8 & 94.2 \\
0.40 & 8.2 & 78.8 \\
\bottomrule
\end{tabularx}
\end{table}


\subsection*{De-identification Pipeline Details}

The de-identification pipeline consists of three sequential stages:
\begin{enumerate}
\item \textbf{Rule-based filtering}: Regular expressions identify and replace structured identifiers including Chinese national ID numbers, phone numbers, email addresses, physical addresses, and hospital record numbers.
\item \textbf{NER-based detection}: A fine-tuned BERT-based NER model identifies contextual identifiers (person names, hospital names, dates, locations) and replaces them with type-consistent placeholders (e.g., \texttt{[\,PERSON\,]}, \texttt{[\,HOSPITAL\,]}, \texttt{[\,DATE\,]}).
\item \textbf{Manual review}: Two independent annotators review all de-identified records. Disagreements are resolved by a third annotator.
 recall and 97.2\% precision.
\end{enumerate}


